\part{The trivially valued theory and envelopes}
\label{part:bj}

Over a trivially valued field, the scaling action makes homogeneity,
filtrations, valuations, and b-divisors central.  The basic formalism is
developed in
\Cref{chap:trivial-analytification,chap:bj-test-configurations,chap:bj-fs-pl-psh,chap:bj-homogeneity-bdivisors}.
The chapters from \Cref{chap:bj-envelope-statement} through
\Cref{chap:bj-envelope-consequences} present the envelope property for
integral projective unibranch varieties, in arbitrary characteristic and for
real numerical classes, with the hypotheses of each step displayed.  The
argument is divided by mechanism: realization, birational reflection, test
ideals, finite descent, alterations, and normalization.  The remaining
chapters develop energy, finite-energy measures, strong topologies, and
stability functionals.

\paragraph{Foundational proof boundary.}
The critical homogeneous Siu-division argument in
\cite[Theorem~6.21]{BJ22} contains an incorrect monotonicity inference.
\Cref{warn:ch35-siu-source-gap} isolates the extra assertion
$\mathrm{(SD)}_k$. The general envelope proof and its dependent
consequences in this part are conditional on that assertion. The strict
division theorem and the independently published smooth envelope cases
are distinguished at their points of use. Compilation and editorial
completion do not discharge this mathematical dependency.

\paragraph{Prerequisites and support.}
\label{front:part4-prerequisites-and-support}
The valuation and analytification language of
\Cref{chap:valuations,chap:berkovich} is required, as are the metric and model
examples in \Cref{chap:metrics,chap:model-calculations}.  The proof of the
envelope property also uses the birational, ideal-theoretic, and b-divisor
material collected in
\Cref{app:birational-geometry,app:ideals,app:intersection-bdivisors,app:stacks}.
Readers interested only in filtrations and finite energy can begin with less
scheme theory and consult those appendices when the envelope proof invokes
them.

\paragraph{Conceptual clusters.}
\label{front:part4-conceptual-clusters}
The valuation and scaling foundation is
\Cref{chap:trivial-analytification,chap:bj-test-configurations,chap:bj-fs-pl-psh,chap:bj-homogeneity-bdivisors}.
It passes from semivaluations and the scaling cone to test configurations,
graded norms, Fubini--Study functions, and b-divisors.  The history and exact
statement of the envelope problem are separated in
\Cref{chap:bj-envelope-history,chap:bj-envelope-statement}.  Its proof runs
through
\Cref{chap:bj-cartier-realization,chap:bj-reflection,chap:bj-smooth-characteristic,chap:bj-finite-descent,chap:bj-alterations,chap:bj-universal-homeomorphisms,chap:bj-subintegral-normalization},
and the deductions are gathered in \Cref{chap:bj-envelope-consequences}.
The variational and metric completion is
\Cref{chap:bj-energy,chap:bj-finite-energy,chap:bj-strong-topologies,chap:bj-linear-growth,chap:bj-stability}.

\paragraph{Reading routes.}
\label{front:part4-reading-routes}
The \emph{short route} follows
\Cref{chap:trivial-analytification,chap:bj-test-configurations,chap:bj-fs-pl-psh,chap:bj-envelope-statement,chap:bj-envelope-consequences,chap:bj-energy};
the mechanism summaries in the proof chapters may be consulted as needed.
The \emph{standard route} reads the homogeneous foundation, then every chapter
from the envelope statement through its consequences, and finishes with
finite energy and strong topology.  The \emph{reference route} reads the part
linearly with
\Cref{app:valuation-algebra,app:birational-geometry,app:ideals,app:intersection-bdivisors}
open beside it; this is the route needed to audit characteristic and descent
hypotheses.  Readers focused on stability should add
\Cref{chap:bj-linear-growth,chap:bj-stability} and then continue to
\Cref{chap:rays-filtrations,chap:ytd-bridge}.

\paragraph{How to read the envelope proof.}
\label{front:part4-how-to-read-the-envelope-proof}
The characteristic-zero route through multiplier ideals gives useful
intuition but is not the arbitrary-characteristic argument.  Alterations and
universal homeomorphisms cannot be omitted when the ground field or the
variety is not already smooth and normal.  Likewise, Cartier realization,
descent across finite maps, and passage through normalization solve different
problems and should not be compressed into a single ``birational invariance''
slogan.

\paragraph{Main outcomes and interfaces.}
\label{front:part4-main-outcomes-and-interfaces}
By the end of the part, the reader should be able to pass between homogeneous
FS functions, graded norms, test configurations, and Cartier b-divisors; to
locate the exact descent input in the envelope argument; and to distinguish
the model, finite-energy, and linear-growth domains of the energy and
stability functionals.  The projective complex interface runs through
filtrations, multiplier ideals, and geodesic rays in
\Cref{chap:valuative-analytification,chap:dxz-bj-comparison,chap:rays-filtrations}.
The BFJ theory shares the model-intersection core, while CLD and
Gubler--Rabinoff supply local differential and curve-test notions; neither
interface identifies the full psh classes without the hypotheses recorded in
\Cref{chap:master-dictionary}.

\paragraph{Normalization and scope.}
\label{front:part4-normalization-and-scope}
The scaling action is retained on the whole trivially valued cone.
BJ psh functions can be inhomogeneous;
\Cref{chap:bj-homogeneity-bdivisors} recovers them from homogeneous
slices through the stated Legendre correspondence. Divisorial valuations are not
silently normalized as they are over a discretely valued field, and points
with nonzero support remain part of the analytification.  Probability
measures are used only after division by a positive volume; zero-volume
classes remain in unnormalized inverse systems.  The envelope theorem has
its stated integral, projective, and unibranch domain, while finite-energy and
stability statements impose their own polarization and linear-growth
hypotheses.
